\section{Introduction}
\label{sec:intro}

\todo{full intro draft.}

\paragraph{Problem.}
Vector-memory-backed multi-agent LLM systems have become the default
architecture for long-horizon tasks (planning, code review, research
assistants). Shared memory is a force multiplier for throughput but it
is also a single high-value attack surface: an adversary who can write
even one malicious document into the shared store can influence every
retrieval-time decision that follows. Recent work (MINJA, AgentPoison,
A-MemGuard) has shown that such attacks are both feasible and
hard to detect with static input filters.

\paragraph{Gap.}
Existing defenses are predominantly reactive: they inspect each
retrieval or each generated action in isolation, after the poison has
already taken effect. Temporal context -- whether an agent's behavior
has drifted across a sequence of tasks -- is largely unused.

\paragraph{Contribution.}
This paper argues that agent-level temporal monitoring is both
sufficient to flag memory-poisoning attacks early and cheap enough to
run inline. Concretely, we contribute:
\begin{enumerate}
  \item \textbf{A six-signal behavioral profile} per agent, measured
  on every PR-handling step (\cref{sec:system}).
  \item \textbf{A BOCPD + Chronos forecaster ensemble} that flags
  signal-level change points before the attacker's payload stabilizes
  into a new equilibrium (\cref{sec:methodology}).
  \item \textbf{A health-weighted decentralized allocator} that routes
  tasks away from drifting agents without pausing the whole system
  (\cref{sec:system}).
  \item \textbf{An end-to-end evaluation} against MINJA and AgentPoison
  with three ablations and two reactive baselines, reporting Advance
  Warning Time (\awt{}), efficiency under attack, and ROC curves
  (\cref{sec:experiments,sec:results}).
\end{enumerate}

\paragraph{Roadmap.}
\Cref{sec:related} surveys memory-poisoning attacks, multi-agent
allocators, and online change-point detection.
\Cref{sec:system} presents the \sys{} architecture.
\Cref{sec:methodology} details the signals, BOCPD, and forecaster.
\Cref{sec:experiments} enumerates the six experimental configurations
and maps them to claims C1--C5.
\Cref{sec:results} reports the measurements.
\Cref{sec:discussion,sec:conclusion} close with limitations and future
work.
